% Context from chapters/shannon-coding.tex; for mathematical reference, not compilation.
\section{Source Coding}
\label{sec-shannon-source}

\newcommand{\Symb}{s}

   
\paragraph{Uniform coding.}

We consider an alphabet $(\Symb_1,\ldots,\Symb_K)$ of $K$ symbols. For signal samples $0 \leq f_0 < 1$ quantized into $K$ bins, the symbols label the bins. For text, they may represent letters, punctuation marks, and spaces.
 
A fixed-length binary code uses $\lceil \log_2(K) \rceil$ bits per symbol.
For example, if $K=4$ and the symbols are $\{0,1,2,3\}$, one possible code is $(c_0=00,c_1=01,c_2=10,c_3=11)$.

When some symbols are much more frequent than others, variable-length codes can reduce the average number of bits per symbol by assigning shorter codewords to frequent symbols.

   
\paragraph{Prefix coding.}
 

A binary code assigns each symbol $\Symb_k$ a finite word $c_k=c(\Symb_k)\in\bigcup_{\ell\geq0}\{0,1\}^{\ell}$ of length $|c_k|$. An empty word is allowed only for a one-symbol alphabet, with the message length supplied separately.
 
A code is a prefix code if no codeword is a prefix of another. Equivalently, the codewords correspond to leaves of a binary tree $T$: starting at the root, traversing an edge to the left appends 0, and traversing an edge to the right appends 1.
 
We write $c=\text{Leaves}(T)$ for this tree representation.


\myfigure{
\image{1-shannon}{.35}{tree/prefix-coding}
\qquad
\image{1-shannon}{.35}{tree/tree-example}
}{ 
Left: the complete tree of binary words of length 3. Right: a prefix code.
}{fig-tree-prefix}

To decode a binary stream, follow its bits down the tree. Upon reaching a leaf, output the corresponding symbol and return to the root. The prefix property makes this decoding unambiguous.


  
\paragraph{Probabilistic modeling.}

We seek a code with the smallest expected length.
 
Let $p=(p_1,\ldots,p_K)\in\RR_+^K$ be the symbol probability vector, with $\sum_kp_k=1$.
 
In practice, one often estimates $p_k$ from the relative frequency of $\Symb_k$ in the data. The decoder must also know the resulting code, either from a shared model or from a transmitted description.

\wrapf{1-shannon/entropie/entropie-fonction}{The concave function $h(u)=-u\log_2u$ underlying entropy.}

The entropy of the probability distribution $p$ is
\eq{
	H(p) \eqdef -\sum_k p_k \log_2(p_k)
}
with the convention $0\log_2(0)=0$. 

For $h(u)=-u\log_2u$ and $u>0$, one has $h'(u)=-(\ln u+1)/\ln2$ and $h''(u)=-1/(u\ln2)<0$. Thus $H$ is strictly concave on the probability simplex.

For a probability density $f$ with respect to a reference measure $\d x$, such as Lebesgue measure on $\RR^d$, differential entropy is defined by $H(f)\eqdef-\int f(x)\log_2(f(x))\,\d x$, whenever this integral is defined. Unlike discrete entropy, it can be negative and depends on the reference measure.

\myfigure{
\tabtrois{
\includegraphics[width=.25\linewidth]{1-shannon/entropie/histo-1}&
\includegraphics[width=.25\linewidth]{1-shannon/entropie/histo-3}&
\includegraphics[width=.25\linewidth]{1-shannon/entropie/histo-sub}\\
$H(p)=0$ & $H(p)=\log_2(4)=2$ &  $H(p) \approx 1.62$
}
}{Three examples of probability distributions with corresponding entropies in bits.}{fig-histo-entropy}


\begin{figure}[tbp]
\centering
\BookDrawing{tikz/coding-entropy-extrema}
\caption{Entropy is minimized at a point mass and maximized at the uniform distribution. Left: level sets of the entropy formula in the positive quadrant and the constraint $p_1+p_2=1$. Middle: its restriction to that constraint. Right: entropy level sets on the three-symbol probability simplex.}
\label{fig-entropy-extrema}
\end{figure}



\begin{lem} For every probability vector $p$,
\eq{
	0 = H(\de_i) \leq H(p) \leq H(1/K) = \log_2(K)
}
where $\de_i$ is the point mass at symbol $i$, and $1/K$ denotes the uniform probability vector.
\end{lem}


\begin{proof}
For $u\in[0,1]$, $-u\log_2u\geq0$, so $H(p)\geq0$, with equality precisely for a point mass.
For probability vectors $p,q$ with $q_i>0$ whenever $p_i>0$, the inequality $\ln u\leq u-1$ gives
\eq{
 \sum_{i:p_i>0}p_i\log_2\frac{q_i}{p_i}
 \leq\frac{1}{\ln2}\left(\sum_{i:p_i>0}q_i-1\right)\leq0.
}
Consequently $H(p)\leq-\sum_i p_i\log_2q_i$. Taking $q_i=1/K$ yields $H(p)\leq\log_2K$, 